\chapter{HASIL DAN PEMBAHASAN}

\section{Karakteristik Dataset}

Dataset final pada penelitian ini diposisikan sebagai artefak pengetahuan operasional, bukan sekadar kumpulan data historis \textit{invoice}. Dataset tersebut merupakan hasil dari proses \textit{Knowledge Acquisition} dan \textit{Knowledge Formalization} yang mengubah regulasi pelanggan, jadwal penerimaan dokumen, batas \textit{cutoff}, dan praktik administrasi menjadi atribut yang dapat digunakan dalam evaluasi representasi pengetahuan. Dengan demikian, karakteristik dataset perlu dibaca sebagai dasar untuk menilai apakah pengetahuan operasional yang telah diformalkan dapat dipertahankan maknanya dalam dua bentuk representasi.

Dataset final yang digunakan telah dibersihkan dan diberi label berdasarkan \textit{Operational Labeling Guideline}. Dataset tersebut berisi 105 data \textit{invoice} dengan 105 nomor \textit{invoice} unik, sehingga tidak terdapat duplikasi nomor \textit{invoice}. Kolom target yang digunakan adalah \texttt{final\_label}, yaitu representasi implementatif dari \textit{Guideline-Based Ground Truth}. Hasil validasi akhir menunjukkan bahwa tidak terdapat label target yang kosong, tidak terdapat nilai label yang tidak valid, serta tidak terdapat nilai \texttt{cutoff\_rule} yang tidak sesuai dengan kategori operasional yang digunakan dalam penelitian.

Distribusi label pada dataset final ditunjukkan pada Tabel \ref{tab:distribusi-label-final}. Komposisi ini penting karena dataset memuat tiga kondisi prioritas operasional, yaitu HIGH, MEDIUM, dan NORMAL. Keberadaan tiga kelas tersebut menunjukkan bahwa pedoman operasional tidak hanya membedakan dokumen mendesak dan tidak mendesak, tetapi juga menangkap kondisi antara yang membutuhkan pemantauan administratif.

\begin{table}[htbp]
    \centering
    \caption{Distribusi Label pada Dataset Final}
    \label{tab:distribusi-label-final}
    \renewcommand{\arraystretch}{1.25}
    \begin{tabular}{lrr}
        \toprule
        \textbf{Kelas Prioritas} & \textbf{Jumlah Data} & \textbf{Persentase} \\
        \midrule
        HIGH & 27 & 25,71\% \\
        MEDIUM & 33 & 31,43\% \\
        NORMAL & 45 & 42,86\% \\
        \midrule
        \textbf{Total} & \textbf{105} & \textbf{100,00\%} \\
        \bottomrule
    \end{tabular}
\end{table}

Karakteristik utama dataset ini terletak pada keberadaan atribut operasional yang merepresentasikan proses administrasi \textit{invoice}, seperti jenis aturan \textit{cutoff}, jarak hari menuju \textit{cutoff}, jadwal penerimaan dokumen, kode hari penerimaan, jarak menuju jadwal penerimaan berikutnya, serta indikator akhir bulan. Atribut-atribut tersebut merupakan hasil formalisasi konsep operasional. Oleh karena itu, dataset ini menjadi penghubung antara pengetahuan yang diperoleh dari praktik administrasi dan dua representasi yang dievaluasi pada bab ini: representasi eksplisit melalui Rule Base dan rekonstruksi statistik melalui \textit{Decision Tree} C4.5.

Pembagian data dilakukan menggunakan \textit{stratified random split} dengan \texttt{test\_size = 0.20} dan \texttt{random\_state = 42}. Pembagian secara stratifikasi digunakan agar proporsi kelas HIGH, MEDIUM, dan NORMAL tetap terwakili pada data latih dan data uji. Rincian pembagian data ditunjukkan pada Tabel \ref{tab:split-final}.

\begin{table}[htbp]
    \centering
    \caption{Distribusi Kelas pada Data Latih dan Data Uji}
    \label{tab:split-final}
    \renewcommand{\arraystretch}{1.25}
    \begin{tabular}{lrr}
        \toprule
        \textbf{Kelas Prioritas} & \textbf{Data Latih} & \textbf{Data Uji} \\
        \midrule
        HIGH & 22 & 5 \\
        MEDIUM & 26 & 7 \\
        NORMAL & 36 & 9 \\
        \midrule
        \textbf{Total} & \textbf{84} & \textbf{21} \\
        \bottomrule
    \end{tabular}
\end{table}

Tabel \ref{tab:split-final} menunjukkan bahwa setiap kelas tetap muncul pada data uji. Hal ini membuat evaluasi konsistensi representasi dapat dilakukan terhadap seluruh kategori prioritas yang digunakan dalam pedoman operasional. Dengan demikian, dataset final berfungsi sebagai artefak pengetahuan yang memuat atribut formal, \textit{Guideline-Based Ground Truth}, dan dasar evaluasi untuk melihat apakah representasi Rule-Based dan rekonstruksi C4.5 mempertahankan makna operasional yang sama.

\section{Representasi Pengetahuan Rule-Based}

Model \textit{Rule-Based} pada penelitian ini diposisikan sebagai representasi eksplisit dari pengetahuan operasional yang telah diformalkan. Rule Base R1 sampai R12 tidak diperlakukan sebagai algoritma pesaing, tetapi sebagai artefak \textit{Knowledge Representation} yang menerjemahkan \textit{Operational Labeling Guideline} ke dalam aturan IF-THEN. Setiap aturan menyimpan makna administratif tertentu, misalnya kedekatan terhadap \textit{cutoff}, keterbatasan jadwal penerimaan, kondisi akhir bulan, atau kondisi tanpa \textit{cutoff}.

Nilai ilmiah dari representasi Rule-Based terletak pada keterlacakan dan transparansi. Ketika sebuah \textit{invoice} diberi rekomendasi prioritas HIGH, MEDIUM, atau NORMAL, alasan keputusan dapat ditelusuri melalui Rule ID yang aktif. Keterlacakan ini penting karena prioritas pengiriman \textit{invoice} berhubungan dengan aturan pelanggan dan konsekuensi administratif. Dengan kata lain, Rule Base mempertahankan hubungan langsung antara pedoman operasional dan rekomendasi prioritas.

Frekuensi aktivasi aturan pada keseluruhan dataset ditunjukkan pada Tabel \ref{tab:rule-usage-final}. Aturan R10 dan R8 menjadi aturan yang paling sering aktif. Pola ini menunjukkan bahwa sebagian besar data berada pada kondisi normal atau mendekati \textit{cutoff} tetapi belum memenuhi kondisi prioritas tertinggi. Aturan R2, R3, dan R4 lebih jarang aktif karena kondisi operasional yang dipenuhi lebih spesifik. Distribusi aktivasi aturan ini memperlihatkan bagaimana pengetahuan operasional yang diformalkan digunakan secara berbeda pada kasus-kasus \textit{invoice}.

\begin{table}[htbp]
    \centering
    \caption{Frekuensi Aktivasi Aturan Rule-Based pada Dataset Final}
    \label{tab:rule-usage-final}
    \renewcommand{\arraystretch}{1.25}
    \begin{tabular}{lr}
        \toprule
        \textbf{Rule ID} & \textbf{Frekuensi} \\
        \midrule
        R2 & 1 \\
        R3 & 3 \\
        R4 & 4 \\
        R6 & 5 \\
        R7 & 17 \\
        R8 & 30 \\
        R9 & 7 \\
        R10 & 38 \\
        R11 & 7 \\
        \bottomrule
    \end{tabular}
\end{table}

Hasil konsistensi Rule-Based terhadap \textit{Guideline-Based Ground Truth} pada data uji ditunjukkan melalui \textit{confusion matrix} pada Gambar \ref{fig:cm-rule-based-final}. Pada data uji sebanyak 21 \textit{invoice}, keluaran Rule-Based berada pada diagonal matriks untuk kelas HIGH, MEDIUM, dan NORMAL.

\begin{figure}[htbp]
    \centering
    \includegraphics[width=0.72\textwidth]{\detokenize{final_experiment/figures/confusion_matrix_rule_based.png}}
    \caption{\textit{Confusion Matrix} Model Rule-Based}
    \label{fig:cm-rule-based-final}
\end{figure}

Nilai akurasi, presisi, \textit{recall}, dan F1-score pada data uji mencapai 1,000 untuk setiap kelas. Dalam posisi ilmiah penelitian ini, angka tersebut dibaca sebagai konsistensi internal antara \textit{Operational Labeling Guideline}, \textit{Guideline-Based Ground Truth}, dan Rule Base. Angka tersebut tidak digunakan untuk menyatakan kebenaran universal keputusan prioritas, tetapi untuk menunjukkan bahwa pengetahuan yang diformalkan dapat direpresentasikan secara eksplisit dan diterapkan secara konsisten pada dataset final.

\section{Rekonstruksi Pengetahuan Decision Tree C4.5}

Model \textit{Decision Tree} C4.5 pada penelitian ini diposisikan sebagai mekanisme \textit{Knowledge Reconstruction}. Model ini tidak menerima Rule Base R1 sampai R12 secara langsung. Masukan yang digunakan adalah atribut operasional yang telah diformalkan dan \textit{Guideline-Based Ground Truth}. Oleh karena itu, hasil C4.5 harus ditafsirkan sebagai rekonstruksi statistik terhadap struktur pengetahuan yang telah diformalkan, bukan sebagai pembentukan aturan operasional di luar pedoman.

Struktur pohon keputusan hasil pelatihan ditunjukkan pada Gambar \ref{fig:decision-tree-final}. Akar pohon menggunakan atribut \texttt{days\_to\_cutoff\_decision\_time}. Posisi atribut ini pada akar pohon menunjukkan dominasi pengetahuan tentang jarak menuju batas administratif. Secara operasional, hal tersebut sejalan dengan pedoman prioritas karena kedekatan terhadap \textit{cutoff} merupakan salah satu konsep utama dalam penentuan urgensi pengiriman \textit{invoice}.

\begin{figure}[htbp]
    \centering
    \includegraphics[width=\textwidth]{\detokenize{final_experiment/figures/decision_tree_visualization.png}}
    \caption{Visualisasi Decision Tree C4.5}
    \label{fig:decision-tree-final}
\end{figure}

Pohon keputusan membentuk beberapa jalur keputusan utama. Apabila \texttt{days\_to\_cutoff\_decision\_time} bernilai kurang dari atau sama dengan 3,5, model mengarah pada kelas HIGH. Pada kondisi lain, struktur pohon mempertimbangkan jadwal penerimaan harian, kode hari penerimaan, dan jarak menuju jadwal penerimaan berikutnya. Jalur-jalur ini menunjukkan bahwa C4.5 merekonstruksi konsep operasional yang sama dengan pedoman, yaitu kombinasi antara kedekatan terhadap \textit{cutoff} dan keterbatasan kesempatan penerimaan dokumen.

Hasil konsistensi C4.5 terhadap \textit{Guideline-Based Ground Truth} pada data uji ditunjukkan pada Gambar \ref{fig:cm-dt-final}. Pola matriks tersebut menunjukkan bahwa keluaran jalur keputusan C4.5 sejalan dengan label berbasis pedoman pada data uji.

\begin{figure}[htbp]
    \centering
    \includegraphics[width=0.72\textwidth]{\detokenize{final_experiment/figures/confusion_matrix_decision_tree.png}}
    \caption{\textit{Confusion Matrix} Model Decision Tree C4.5}
    \label{fig:cm-dt-final}
\end{figure}

Interpretasi utama dari hasil C4.5 bukan terletak pada perayaan metrik, melainkan pada struktur rekonstruksi. Model membentuk jalur yang dapat dibaca dan dibandingkan dengan Rule Base. Karena C4.5 tidak menerima aturan eksplisit, kesesuaian jalur pohon dengan pedoman menunjukkan bahwa atribut operasional yang diformalkan memuat struktur pengetahuan yang dapat direkonstruksi secara statistik.

Peringkat kepentingan fitur pada Gambar \ref{fig:feature-importance-final} dan Tabel \ref{tab:feature-importance-final} digunakan untuk melihat dominasi pengetahuan dalam proses rekonstruksi. Kepentingan fitur tidak ditafsirkan sebagai bukti kausal, tetapi sebagai indikasi atribut formal mana yang paling banyak digunakan oleh pohon untuk mempertahankan struktur keputusan berbasis pedoman.

\begin{figure}[htbp]
    \centering
    \includegraphics[width=0.82\textwidth]{\detokenize{final_experiment/figures/feature_importance.png}}
    \caption{Kepentingan Fitur pada Decision Tree C4.5}
    \label{fig:feature-importance-final}
\end{figure}

\begin{table}[htbp]
    \centering
    \caption{Peringkat Kepentingan Fitur Decision Tree C4.5}
    \label{tab:feature-importance-final}
    \renewcommand{\arraystretch}{1.25}
    \begin{tabular}{lr}
        \toprule
        \textbf{Fitur} & \textbf{Importance} \\
        \midrule
        \texttt{days\_to\_cutoff\_decision\_time} & 0,4818 \\
        \texttt{receive\_day\_code} & 0,2984 \\
        \texttt{next\_receive\_day\_gap\_feature} & 0,1429 \\
        \texttt{receive\_schedule} & 0,0769 \\
        Fitur operasional lain & 0,0000 \\
        \bottomrule
    \end{tabular}
\end{table}

Tabel \ref{tab:feature-importance-final} memperlihatkan bahwa fitur non-nol dalam rekonstruksi terdiri atas \texttt{days\_to\_cutoff\_decision\_time}, \texttt{receive\_day\_code}, \texttt{next\_receive\_day\_gap\_feature}, dan \texttt{receive\_schedule}. Keempat fitur tersebut dapat dibaca sebagai tiga kelompok pengetahuan operasional dominan: jarak terhadap \textit{cutoff}, informasi hari penerimaan, dan batasan jadwal penerimaan. Dominasi ini selaras dengan proses formalisasi karena ketiga kelompok tersebut merupakan dasar operasional dalam penentuan prioritas pengiriman.

Sebaliknya, fitur seperti \texttt{cutoff\_rule}, \texttt{receive\_month\_end\_flag}, \texttt{receive\_weekday\_code}, dan \texttt{receive\_day\_of\_month} memiliki nilai kepentingan 0 pada model akhir. Kondisi ini tidak berarti bahwa fitur-fitur tersebut tidak relevan secara administrasi. Dalam struktur pohon yang terbentuk, sebagian informasi dari fitur tersebut telah direpresentasikan secara lebih langsung oleh \texttt{days\_to\_cutoff\_decision\_time} dan atribut jadwal penerimaan. Dengan demikian, nilai kepentingan 0 dibaca sebagai redundansi dalam rekonstruksi pohon, bukan sebagai ketiadaan makna operasional.

\section{Analisis Representasi Komparatif}

Analisis komparatif pada penelitian ini tidak digunakan untuk menentukan peringkat paradigma. Analisis ini digunakan untuk menjelaskan perbedaan mekanisme representasi pengetahuan. Rule-Based merepresentasikan pengetahuan operasional secara eksplisit melalui aturan, sedangkan C4.5 merekonstruksi pengetahuan yang sama dari atribut operasional dan \textit{Guideline-Based Ground Truth}. Dengan demikian, perbandingan diarahkan pada sumber pengetahuan, transparansi, interpretabilitas, pemeliharaan, kemampuan rekonstruksi, dan kesesuaian makna operasional.

Perbandingan metrik evaluasi ditunjukkan pada Tabel \ref{tab:metric-comparison-final}. Tabel tersebut digunakan sebagai bukti pendukung bahwa kedua representasi menghasilkan keluaran yang konsisten terhadap \textit{Guideline-Based Ground Truth} pada data uji. Metrik tidak digunakan sebagai dasar pemeringkatan antarmodel.

\begin{table}[htbp]
    \centering
    \caption{Perbandingan Metrik Evaluasi pada Data Uji}
    \label{tab:metric-comparison-final}
    \scriptsize
    \renewcommand{\arraystretch}{1.25}
    \begin{tabular}{lrrrrr}
        \toprule
        \textbf{Model} & \textbf{Akurasi} & \textbf{Prec. Makro} & \textbf{Rec. Makro} & \textbf{Macro F1} & \textbf{Weighted F1} \\
        \midrule
        Rule-Based & 1,000 & 1,000 & 1,000 & 1,000 & 1,000 \\
        Decision Tree C4.5 & 1,000 & 1,000 & 1,000 & 1,000 & 1,000 \\
        \bottomrule
    \end{tabular}
\end{table}

Perbedaan konseptual antara kedua pendekatan dirangkum pada Tabel \ref{tab:comparison-characteristics}. Tabel ini menempatkan \textit{Rule-Based} sebagai representasi pengetahuan eksplisit dan C4.5 sebagai rekonstruksi pola keputusan berbasis data. Dengan demikian, pembacaan tabel harus dilakukan pada level paradigma representasi, bukan pada level kompetisi performa.

\begin{table}[htbp]
    \centering
    \caption{Perbandingan Karakteristik Rule-Based dan Decision Tree C4.5}
    \label{tab:comparison-characteristics}
    \scriptsize
    \renewcommand{\arraystretch}{1.25}
    \begin{tabularx}{\textwidth}{p{0.21\textwidth}XX}
        \toprule
        \textbf{Aspek} & \textbf{Rule-Based} & \textbf{Decision Tree C4.5} \\
        \midrule
        Sumber pengetahuan & Berasal dari \textit{Operational Labeling Guideline} yang diformalkan menjadi aturan R1 sampai R12. & Berasal dari data latih yang memuat atribut operasional dan label akhir berbasis pedoman. \\
        Mekanisme keputusan & Menggunakan pencocokan kondisi IF-THEN terhadap aturan yang telah didefinisikan. & Menggunakan pemisahan data berdasarkan \textit{entropy} dan menghasilkan jalur keputusan dari akar ke daun. \\
        Interpretabilitas & Interpretasi langsung melalui Rule ID dan kondisi aturan yang aktif. & Interpretasi dilakukan melalui jalur pohon dan nilai fitur pada setiap node. \\
        Transparansi & Sangat eksplisit karena struktur keputusan ditentukan sebelum prediksi dilakukan. & Transparan setelah pohon dibentuk, tetapi struktur keputusannya dipengaruhi oleh data latih. \\
        Maintainability & Perubahan regulasi dapat diterjemahkan dengan memperbarui aturan tertentu. & Perubahan pola operasional memerlukan data baru dan pelatihan ulang model. \\
        Adaptabilitas & Stabil untuk lingkungan dengan aturan administratif yang tetap. & Memerlukan data berlabel baru untuk merepresentasikan perubahan operasional. \\
        Kebutuhan pelatihan & Tidak memerlukan proses pelatihan statistik. & Memerlukan data latih, proses encoding, dan pembentukan model. \\
        Representasi pengetahuan & Pengetahuan direpresentasikan secara deklaratif dan eksplisit. & Pengetahuan direpresentasikan secara terkompresi dalam struktur pohon. \\
        Kesesuaian operasional & Sesuai untuk kebutuhan audit, penelusuran alasan keputusan, dan dokumentasi aturan administratif. & Sesuai untuk analisis pola, dominasi atribut, dan rekonstruksi pedoman dari data. \\
        Skalabilitas masa depan & Bergantung pada kemampuan menjaga dan memperbarui Rule Base ketika regulasi berubah. & Bergantung pada ketersediaan data yang representatif dan pengendalian risiko overfitting. \\
        \bottomrule
    \end{tabularx}
\end{table}

Tabel \ref{tab:comparison-characteristics} menunjukkan bahwa kedua pendekatan memiliki fungsi ilmiah yang berbeda. Rule-Based mempertahankan kedekatan dengan bahasa administratif karena aturan dapat dibaca sebagai prosedur operasional. C4.5 memberikan struktur data-driven yang dapat diperiksa melalui jalur pohon dan dominasi fitur. Keduanya menjawab pertanyaan berbeda: apakah pengetahuan dapat direpresentasikan secara eksplisit, dan apakah pengetahuan yang sama dapat direkonstruksi dari atribut operasional.

Dari sisi interpretabilitas, Rule-Based menjelaskan keluaran melalui Rule ID dan kondisi aturan yang aktif, sedangkan C4.5 menjelaskan keluaran melalui jalur dari akar pohon ke daun. Penjelasan Rule-Based lebih dekat dengan bahasa pedoman operasional. Penjelasan C4.5 lebih dekat dengan struktur atribut karena menggunakan ambang seperti 3,5 dan 10,5 hari sebagai hasil pemisahan data.

Dari sisi pemeliharaan, Rule-Based dapat diperbarui melalui perubahan aturan ketika regulasi pelanggan berubah. C4.5 memerlukan data baru dan pelatihan ulang agar struktur pohon mencerminkan kondisi baru. Perbedaan ini menunjukkan bahwa kedua pendekatan memiliki peran yang saling melengkapi dalam kerangka formalisasi pengetahuan: satu menjaga transparansi aturan, sedangkan yang lain menyediakan rekonstruksi dan analisis dominasi atribut.

Ringkasan kesalahan klasifikasi pada data uji ditunjukkan pada Tabel \ref{tab:misclassification-summary-final}. Tidak adanya \textit{false positive}, \textit{false negative}, maupun kesalahan antar kelas dibaca sebagai konsistensi representasi terhadap \textit{Guideline-Based Ground Truth} pada dataset final. Temuan ini tidak digunakan untuk memperluas klaim ke luar konteks penelitian, melainkan sebagai bukti bahwa kedua representasi mempertahankan makna operasional yang sama dalam ruang data yang dievaluasi.

\begin{table}[htbp]
    \centering
    \caption{Ringkasan Kesalahan Klasifikasi pada Data Uji}
    \label{tab:misclassification-summary-final}
    \renewcommand{\arraystretch}{1.25}
    \begin{tabular}{lrrr}
        \toprule
        \textbf{Model} & \textbf{False Positive} & \textbf{False Negative} & \textbf{Total Kesalahan} \\
        \midrule
        Rule-Based & 0 & 0 & 0 \\
        Decision Tree C4.5 & 0 & 0 & 0 \\
        \bottomrule
    \end{tabular}
\end{table}

\section{\textit{Knowledge Representation Perspective}}

Subbab ini merupakan bagian utama dari Chapter IV karena menghubungkan hasil eksperimen dengan kontribusi ilmiah penelitian. Berdasarkan \textit{Research Positioning Statement}, kontribusi utama penelitian ini adalah \textit{Operational Knowledge Formalization Framework}. Oleh karena itu, hasil pada bab ini harus dibaca sebagai evaluasi konsistensi representasi pengetahuan, bukan sebagai pemeringkatan prediktif suatu model.

Kebaruan yang ditunjukkan pada subbab ini bukan terletak pada penggunaan Rule-Based atau C4.5 secara terpisah, melainkan pada proses transformasi pengetahuan operasional menjadi dua representasi yang dapat diperiksa konsistensinya. Dengan demikian, pembaca dapat menelusuri bagaimana pengetahuan administratif yang semula berada pada regulasi pelanggan dan praktik kerja berubah menjadi atribut formal, Rule Base, rekonstruksi pohon, dan dasar rekomendasi prioritas.

Alur pengetahuan dalam penelitian ini dimulai dari \textit{Knowledge Acquisition}. Pada tahap tersebut, pengetahuan administratif mengenai batas \textit{cutoff}, jadwal penerimaan dokumen, kondisi prioritas, dan praktik penanganan \textit{invoice} diperoleh dari konteks operasional. Pengetahuan tersebut kemudian melalui \textit{Knowledge Formalization}, yaitu proses menerjemahkan konsep operasional menjadi atribut formal dan aturan eksplisit. Hasil formalisasi ini menghasilkan dua cabang representasi: Rule Base sebagai representasi eksplisit dan dataset operasional sebagai dasar rekonstruksi C4.5.

\textit{Operational Labeling Guideline} berperan sebagai penghubung antara aturan operasional dan \textit{Guideline-Based Ground Truth}. Label pada dataset final tidak diposisikan sebagai kebenaran universal yang berdiri sendiri, tetapi sebagai keputusan referensi yang diturunkan dari pedoman operasional yang telah diformalkan. Oleh karena itu, evaluasi pada bab ini menilai apakah Rule-Based dan C4.5 mempertahankan makna keputusan yang sama terhadap referensi tersebut.

Dalam representasi Rule-Based, pengetahuan digunakan secara langsung. Setiap aturan IF-THEN merepresentasikan kondisi administratif tertentu dan menghasilkan prioritas yang sesuai dengan pedoman. Dalam rekonstruksi C4.5, pengetahuan tidak diberikan sebagai daftar aturan. Model hanya menerima atribut operasional dan \textit{Guideline-Based Ground Truth}. Struktur pohon yang terbentuk karena itu menggambarkan rekonstruksi statistik terhadap pedoman operasional, bukan pembentukan aturan baru di luar proses formalisasi.

\subsection{\textit{Knowledge Compression}}

Kompresi pengetahuan merupakan temuan penting dalam pembacaan struktur C4.5. Rule Base mendeskripsikan pengetahuan operasional melalui 12 aturan eksplisit, sedangkan struktur C4.5 merangkum pola tersebut menjadi tujuh jalur keputusan utama sebagaimana ditunjukkan pada Tabel \ref{tab:mapping-dt-rulebase}. Kompresi ini terjadi karena beberapa aturan memiliki efek keputusan yang serupa ketika diproyeksikan ke dalam atribut formal, terutama atribut yang berkaitan dengan kedekatan terhadap \textit{cutoff} dan keterbatasan jadwal penerimaan.

Kompresi tersebut tidak berarti bahwa jalur pohon menggantikan Rule Base. Rule Base tetap mempertahankan makna administratif secara deklaratif dan dapat ditelusuri ke pedoman operasional. C4.5 menyajikan bentuk terkompresi dari pola yang sama, sehingga beberapa kondisi administratif yang berbeda dapat muncul sebagai satu jalur pohon apabila menghasilkan pemisahan kelas yang serupa pada dataset final. Dengan demikian, kompresi pengetahuan menunjukkan perubahan bentuk representasi, bukan pengurangan makna operasional.

Berdasarkan Tabel \ref{tab:feature-importance-final}, rekonstruksi C4.5 didominasi oleh atribut \texttt{days\_to\_cutoff\_decision\_time}, \texttt{receive\_day\_code}, \texttt{next\_receive\_day\_gap\_feature}, dan \texttt{receive\_schedule}. Keempat fitur non-nol tersebut dapat dikelompokkan menjadi tiga konsep pengetahuan operasional dominan: jarak terhadap \textit{cutoff}, informasi hari penerimaan, dan keterbatasan jadwal penerimaan. Hal ini memperlihatkan bahwa kompresi C4.5 tetap berada dalam ruang pengetahuan yang telah diformalkan.

Dari sudut pandang rekayasa pengetahuan, kompresi tersebut dapat dipahami sebagai penyederhanaan representasi dan abstraksi pengetahuan. Beberapa aturan yang berbeda secara administratif dapat memiliki makna keputusan yang sama ketika dilihat melalui atribut dominan. Misalnya, kondisi yang semula dibedakan melalui jenis \textit{cutoff}, hari penerimaan, atau indikator akhir bulan dapat direpresentasikan kembali melalui jarak menuju \textit{cutoff} dan keterbatasan jadwal penerimaan. Oleh karena itu, atribut dengan nilai kepentingan 0 tidak langsung menunjukkan ketiadaan makna operasional, melainkan dapat menunjukkan redundansi atribut dalam struktur pohon yang telah memilih representasi paling langsung.

\subsection{Pemetaan Jalur Rekonstruksi terhadap Rule Base}

Pemetaan jalur keputusan C4.5 terhadap Rule Base ditunjukkan pada Tabel \ref{tab:mapping-dt-rulebase}. Pemetaan ini menggunakan jalur yang tampak pada struktur pohon dan membandingkannya dengan aturan R1 sampai R12 yang telah diformalkan. Tujuannya bukan mencari aturan baru, melainkan memeriksa apakah jalur rekonstruksi memiliki padanan makna dengan Rule Base.

\begin{table}[htbp]
    \centering
    \caption{Pemetaan Jalur Decision Tree terhadap Rule Base}
    \label{tab:mapping-dt-rulebase}
    \scriptsize
    \renewcommand{\arraystretch}{1.25}
    \begin{tabularx}{\textwidth}{p{0.10\textwidth}X p{0.18\textwidth} p{0.18\textwidth}}
        \toprule
        \textbf{Jalur DT} & \textbf{Kondisi Utama} & \textbf{Padanan Rule-Based} & \textbf{Jenis Kesesuaian} \\
        \midrule
        DTR1 & \texttt{days\_to\_cutoff\_decision\_time} $\leq$ 3,5 menghasilkan HIGH & R2, R3, R4, R6 & Partial Match \\
        DTR2 & 3,5 $<$ \texttt{days\_to\_cutoff\_decision\_time} $\leq$ 10,5, jadwal tidak setiap hari, dan \texttt{next\_receive\_day\_gap\_feature} $\leq$ 2 menghasilkan HIGH & R7 & Partial Match \\
        DTR3 & 3,5 $<$ \texttt{days\_to\_cutoff\_decision\_time} $\leq$ 10,5, jadwal tidak setiap hari, dan \texttt{next\_receive\_day\_gap\_feature} $>$ 2 menghasilkan MEDIUM & R8, R9 & Partial Match \\
        DTR4 & 3,5 $<$ \texttt{days\_to\_cutoff\_decision\_time} $\leq$ 10,5 dan jadwal setiap hari menghasilkan MEDIUM & R8 & Exact Match \\
        DTR5 & \texttt{days\_to\_cutoff\_decision\_time} $>$ 10,5, kode hari penerimaan terbatas, dan \texttt{next\_receive\_day\_gap\_feature} $\leq$ 1,5 menghasilkan HIGH & R7 & Exact Match \\
        DTR6 & \texttt{days\_to\_cutoff\_decision\_time} $>$ 10,5, kode hari penerimaan terbatas, dan \texttt{next\_receive\_day\_gap\_feature} $>$ 1,5 menghasilkan MEDIUM & R9 & Exact Match \\
        DTR7 & \texttt{days\_to\_cutoff\_decision\_time} $>$ 10,5 dan kode hari penerimaan normal menghasilkan NORMAL & R10, R11 & Partial Match \\
        \bottomrule
    \end{tabularx}
\end{table}

Tabel \ref{tab:mapping-dt-rulebase} menunjukkan bahwa jalur DTR1 menggabungkan beberapa kondisi HIGH yang berkaitan dengan kedekatan terhadap \textit{cutoff}. Jalur DTR2 dan DTR5 menunjukkan rekonstruksi terhadap kondisi jadwal penerimaan terbatas yang berhubungan dengan R7. Jalur DTR3, DTR4, dan DTR6 menunjukkan bagaimana kondisi MEDIUM direkonstruksi melalui kombinasi jarak terhadap \textit{cutoff} dan kesempatan penerimaan berikutnya. Jalur DTR7 menunjukkan representasi kondisi NORMAL ketika jarak terhadap \textit{cutoff} masih panjang dan kode hari penerimaan berada pada kondisi normal.

Jenis kesesuaian \textit{Exact Match} dan \textit{Partial Match} penting untuk dibaca secara hati-hati. \textit{Exact Match} menunjukkan bahwa jalur pohon memiliki padanan yang dekat dengan satu aturan operasional. \textit{Partial Match} menunjukkan bahwa satu jalur pohon merangkum lebih dari satu kondisi operasional atau hanya menangkap sebagian dari struktur aturan. Kedua jenis kesesuaian tersebut tetap mendukung interpretasi rekonstruksi karena C4.5 memang menghasilkan struktur yang lebih ringkas daripada Rule Base.

\subsection{Konsistensi Representasi Pengetahuan}

Konsistensi representasi pengetahuan terjadi ketika Rule-Based dan C4.5 menghasilkan keluaran yang sama terhadap \textit{Guideline-Based Ground Truth}, sekaligus mempertahankan makna operasional yang dapat dijelaskan. Pada Rule-Based, makna tersebut dipertahankan melalui Rule ID dan kondisi IF-THEN. Pada C4.5, makna tersebut dipertahankan melalui jalur pohon, ambang atribut, dan dominasi fitur operasional.

Tidak adanya kesalahan klasifikasi pada data uji memperkuat interpretasi konsistensi representasi. Namun, karena \textit{Guideline-Based Ground Truth} dibentuk dari pedoman operasional, hasil ini tidak boleh ditafsirkan sebagai pembuktian di luar konteks bahwa keputusan prioritas akan selalu sama pada semua kondisi nyata. Hasil tersebut menunjukkan bahwa dalam dataset final, dua representasi yang berbeda mempertahankan makna keputusan yang sama terhadap pedoman yang sama.

\section{Diskusi}

Gambar~\ref{fig:knowledge_reconstruction} menjadi ringkasan konseptual untuk membaca makna ilmiah dari hasil Chapter IV.

\begin{figure}[htbp]
\centering
\includegraphics[width=\columnwidth]{KnowledgeReconstruction.png}
\caption{Conceptual Illustration of Knowledge Reconstruction and Representation Consistency.}
\label{fig:knowledge_reconstruction}
\end{figure}

Gambar~\ref{fig:knowledge_reconstruction} memperlihatkan bahwa kontribusi penelitian tidak berhenti pada keluaran numerik, tetapi pada hubungan antara \textit{explicit operational knowledge}, \textit{Rule-Based Representation}, \textit{Decision Tree Reconstruction}, dan \textit{Guideline-Based Ground Truth}. Pengetahuan operasional yang telah diformalkan berperan sebagai sumber makna bersama. Dari sumber tersebut, Rule-Based menyatakan pengetahuan secara eksplisit melalui aturan, sedangkan Decision Tree merekonstruksi pola keputusan dari atribut operasional dan label referensi berbasis pedoman.

Kedua bentuk tersebut bersifat \textit{complementary representations}. Rule-Based menjaga keterbacaan administratif karena setiap \textit{Priority Recommendation} dapat ditelusuri ke kondisi operasional yang diformalkan. Decision Tree Reconstruction memberikan sudut pandang rekonstruktif karena struktur pohon memperlihatkan bagaimana atribut seperti jarak terhadap \textit{cutoff} dan jadwal penerimaan cukup untuk membentuk kembali makna keputusan yang sama. Dengan demikian, C4.5 tidak diposisikan sebagai pengganti pedoman, melainkan sebagai mekanisme untuk memeriksa apakah struktur pengetahuan dapat direkonstruksi dari data yang telah diberi \textit{Guideline-Based Ground Truth}.

Makna utama dari gambar tersebut adalah \textit{representation consistency}. Konsistensi tidak hanya berarti kesamaan keluaran, tetapi juga kesetaraan makna operasional antara representasi eksplisit dan rekonstruksi data-driven. Apabila dua representasi menghasilkan keputusan yang sama terhadap pedoman yang sama dan masih dapat dijelaskan melalui konsep operasional yang sama, maka keduanya mempertahankan \textit{equivalent operational meaning}.

Implikasi ilmiahnya adalah bahwa Knowledge Engineering menjadi pusat penelitian. Tanpa akuisisi, strukturisasi, dan formalisasi pengetahuan, atribut operasional hanya menjadi variabel data yang terpisah dari konteks administrasi. Setelah diformalkan, atribut dan aturan tersebut menjadi artefak pengetahuan yang dapat direpresentasikan, direkonstruksi, dan dievaluasi konsistensinya.

Implikasi praktisnya adalah bahwa hasil evaluasi dapat ditempatkan kembali dalam sistem pelacakan \textit{invoice} dan POD. Rekomendasi prioritas tidak berdiri sebagai label teknis semata, tetapi sebagai bagian dari alur informasi operasional yang mendukung pelacakan dokumen, penyimpanan POD, riwayat pengiriman, dan pengambilan keputusan administratif. Dengan demikian, gambar tersebut menegaskan bahwa nilai ilmiah penelitian terletak pada kemampuan mempertahankan makna pengetahuan operasional ketika pengetahuan tersebut bergerak dari pedoman eksplisit menuju rekonstruksi dan konteks sistem informasi.

\section{Threats to Validity}

\subsection{Validitas Internal}

Ancaman terhadap validitas internal terutama berkaitan dengan hubungan antara \textit{Guideline-Based Ground Truth}, \textit{Operational Labeling Guideline}, dan Rule Base. Label akhir dibentuk menggunakan pedoman operasional, sedangkan Rule-Based juga dibangun dari pengetahuan yang sama. Oleh karena itu, kesesuaian hasil Rule-Based terhadap label lebih tepat dipahami sebagai konsistensi implementasi pengetahuan, bukan sebagai pembuktian independen terhadap kebenaran keputusan operasional.

Pada C4.5, ancaman validitas internal muncul karena model belajar dari dataset yang telah diberi label berbasis pedoman. Kesesuaian C4.5 terhadap label menunjukkan rekonstruksi terhadap struktur pedoman pada data uji, bukan validasi bahwa model akan menghasilkan pola yang sama pada seluruh situasi nyata. Selain itu, struktur pohon dipengaruhi oleh pembagian data latih dan data uji, sehingga hasil pada satu pembagian data tetap perlu dibaca secara hati-hati.

\subsection{Validitas Eksternal}

Validitas eksternal berkaitan dengan sejauh mana temuan dapat diterapkan di luar konteks penelitian. Dataset yang digunakan berasal dari karakteristik administrasi \textit{invoice} dan regulasi pelanggan pada organisasi yang menjadi objek penelitian. Apabila organisasi lain memiliki aturan \textit{cutoff}, pola penerimaan dokumen, kebijakan POD, atau prosedur prioritas yang berbeda, Rule Base, atribut operasional, dan struktur pohon keputusan dapat memerlukan penyesuaian.

Dengan demikian, hasil penelitian tidak dapat langsung digeneralisasi ke seluruh proses administrasi \textit{invoice}. Generalisasi hanya dapat dilakukan setelah dilakukan akuisisi pengetahuan baru, formalisasi ulang terhadap aturan organisasi lain, serta evaluasi konsistensi representasi pada dataset dan pedoman yang sesuai dengan konteks tersebut.

\subsection{Validitas Konstruk}

Validitas konstruk berkaitan dengan apakah atribut yang digunakan benar-benar merepresentasikan konsep prioritas operasional. Penelitian ini menggunakan atribut seperti jarak menuju \textit{cutoff}, jadwal penerimaan, kode hari penerimaan, dan jarak menuju jadwal penerimaan berikutnya. Atribut tersebut merepresentasikan aspek penting dari pedoman operasional, tetapi tidak menutup kemungkinan adanya faktor lain yang belum dimodelkan, misalnya kondisi komunikasi pelanggan, perubahan kebijakan mendadak, kendala pengiriman fisik, atau pertimbangan administratif yang tidak tercatat dalam dataset.

Nilai kepentingan fitur juga harus dibaca dalam batas validitas konstruk. Fitur yang memiliki nilai 0 pada pohon akhir tidak otomatis kehilangan makna operasional. Fitur tersebut dapat menjadi redundan karena informasi yang sama telah diserap oleh atribut lain yang lebih langsung, terutama \texttt{days\_to\_cutoff\_decision\_time} dan atribut jadwal penerimaan.

\subsection{Keterbatasan Ruang Lingkup Pengetahuan}

Pengetahuan yang diformalkan dalam penelitian ini dibatasi pada regulasi pelanggan, atribut waktu administrasi, jadwal penerimaan dokumen, dan kondisi prioritas yang tersedia pada konteks organisasi yang diteliti. Ruang lingkup pengetahuan tersebut belum mencakup seluruh kemungkinan faktor operasional, seperti negosiasi khusus dengan pelanggan, perubahan prosedur mendadak, kendala logistik fisik, atau kebijakan internal yang tidak terdokumentasi dalam dataset.

Keterbatasan ini penting karena \textit{Operational Knowledge Formalization Framework} dapat dipindahkan sebagai metode, tetapi isi pengetahuannya tetap bergantung pada konteks organisasi. Dengan demikian, Rule Base dan rekonstruksi C4.5 pada penelitian ini harus dibaca sebagai representasi pengetahuan yang terikat pada ruang operasional yang telah diakuisisi dan diformalkan.

\subsection{Keterbatasan Dataset}

Dataset final berisi 105 data \textit{invoice}. Jumlah ini cukup untuk menunjukkan proses evaluasi konsistensi antara \textit{Guideline-Based Ground Truth}, Rule Base, dan rekonstruksi C4.5 dalam konteks penelitian ini, tetapi masih terbatas untuk menarik kesimpulan yang luas. Variasi pelanggan, variasi aturan, variasi periode waktu, dan perubahan kondisi operasional dapat menghasilkan struktur pengetahuan yang berbeda.

Laporan validasi juga menunjukkan adanya peringatan terkait nilai \textit{schedule gap} yang hilang pada beberapa baris dengan \texttt{receive\_day\_code} tidak tersedia dan nilai \texttt{cutoff\_value} yang hilang pada sebagian regulasi pelanggan. Peringatan tersebut telah ditangani dalam eksperimen, tetapi tetap menunjukkan bahwa kualitas rekonstruksi bergantung pada kelengkapan atribut operasional yang tersedia.

\subsection{Keterbatasan Guideline-Based Labeling}

Label akhir pada eksperimen final tidak berasal dari validasi independen baru pada tahap evaluasi, melainkan dari \textit{Guideline-Based Labeling} berdasarkan \textit{Operational Labeling Guideline}. Pendekatan ini sesuai dengan tujuan penelitian, yaitu mengevaluasi konsistensi antara representasi eksplisit dan rekonstruksi data-driven dari pengetahuan operasional yang sama. Namun, keterbatasannya adalah hasil evaluasi tidak boleh ditafsirkan sebagai bukti bahwa label tersebut selalu identik dengan keputusan manusia pada seluruh situasi nyata.

Dengan kata lain, \textit{Guideline-Based Ground Truth} adalah referensi operasional yang dibangun dari pedoman, bukan kebenaran eksternal yang bebas dari proses formalisasi. Hal ini tidak melemahkan tujuan penelitian, tetapi membatasi jenis klaim yang boleh dibuat: penelitian ini mengklaim konsistensi representasi pengetahuan dalam konteks yang diformalkan, bukan generalisasi keputusan prioritas di luar konteks tersebut.

\subsection{Keterbatasan Generalisasi dan Future Work}

Generalisasi hasil penelitian perlu dilakukan secara bertahap. Penelitian selanjutnya dapat memperluas dataset dengan periode waktu yang lebih panjang, menambahkan pelanggan dengan variasi regulasi yang lebih beragam, serta mengevaluasi model pada data operasional baru. Selain itu, perubahan regulasi pelanggan perlu diamati secara berkala agar Rule Base dan model data-driven tetap merepresentasikan kondisi administrasi yang aktual.

Future work juga perlu memperkuat validasi eksternal terhadap pedoman. Pengujian pada organisasi lain atau periode operasional baru harus dimulai dari akuisisi pengetahuan ulang, bukan hanya menerapkan model yang sama. Dengan cara tersebut, \textit{Operational Knowledge Formalization Framework} dapat diuji sebagai kerangka yang dapat dipindahkan secara metodologis, sementara isi pengetahuannya tetap disesuaikan dengan batas operasional masing-masing konteks.
